\clearpage
\pagestyle{plain}
\begingroup
\setlength{\parindent}{0pt}
\setlength{\parskip}{6pt}
\fontfamily{\familydefault}\selectfont
\begin{center}
\headline{\color{gray}Cover Letter}
\end{center}
\begin{flushright}
\small
\today\\
% 36 Poureslami Avenue\\
% Komeil Street\\
% Tehran 1353764371\\
% Iran
\normalfont
\end{flushright}
% \vspace{1em}
\begin{flushleft}
\small
St. Anna Kinderkrebsforschung GmbH\\
Cresswell Group\\
Zimmermannplatz 10\\
1090 Vienna \\
Austria
% \today\\
\normalsize
\end{flushleft}
% \vspace{1em}
% \noindent\textbf{Subject:} Motivation Letter for "PhD Student (f/m/d) - Computational Cancer Genomics / Bioinformatics"

\noindent Dear Dr. George D. Cresswell, % and the Selection Committee,
\vspace{1em}

% opening
% What makes this position attractive to you? 
% What relevant papers interest you the most? 
% What qualifications/experiences make you a good fit? Please provide examples where you demonstrated your skills

% I am writing to apply for the PhD position in Computational Cancer Genomics/ Bioinformatics in your group at the St. Anna Children's Cancer Research Institute. I am enormously interested in your research activities, and my background completely matches yours. It will be a great opportunity for me to be part of such an outstanding group. So, I would appreciate it if you could possibly consider my application.

% I have completed my master's degree in Numerical Analysis at the Institute for Advanced Studies in Basic Sciences. During this program, I gained extensive knowledge about biophysical modeling through the winter school mentioned in my resume and my own research on combining different imaging techniques to present a 3D model for breast cancer. Considering the resources I had, I narrowed down my efforts to the importance of applying theoretical models to practical use in clinical settings; I focused my thesis mainly on using neural networks as a means to solve partial differential equations. Later, I shifted my gaze towards a comprehensive study of cancer models, especially "Spatio-Genetic and phenotypic modelling elucidates resistance and re-sensitisation to treatment in heterogeneous melanoma", which presents a state-of-the-art mathematical model that mainly focuses on drug resistance to melanoma cancer treatment with focus on similar concerns as to what I observed in your advertisement, namely, concentration on mutational processes and evolution of pediatric cancers to determine which children are most at risk of a poor response to treatment. 

% Therefore, I began skimming through your research. I have found your work, in particular "Model-based tumor subclonal reconstruction" and "Exploiting evolutionary steering to induce collateral drug sensitivity in cancer", quite fascinating. I agree with you that PINNs(Physics-informed neural nets) can produce better results and less unwanted artifacts compared to having the machine find relevance in a data-driven way. Furthermore, I like the way you applied tumour evolution, and also impressed by the way you quantified the uncertainty of your data. On top of that, your methodology and process is quite apealing to me, for example the way you have done your expriments setup when you were studying the resistence of cells to the drug, and your conclusion which not only focuses on the surviving population, but also the cells that died during the experiment, hypothesising the relation of these to the evolutionary hystory of the cell population.

% I believe my experience with data analysis, machine learning, and my strong background in Applied Mathematics, especially in modeling cancer biology, make me a strong candidate. In addition, I gained invaluable experience in teamwork, working under stress and tight deadlines, and team coordination throughout my career as a professional software developer, which requires an absolute mentality of getting things done. As evident from my resume, I was quite active in attending seminars and conferences, which demonstrates my eagerness to engage in academic pursuits. What is not evident from my resume, however, is that my resolution to study cancer’s biology in particular stems from the fact that I lost my aunt to breast cancer. 

% I would like to bring my passion and experience to your research group, to translate my research into clinically useful metrics to develop computational methods that predict which children are most at risk of poor treatment response.

% In general, my research interests include oncology, biophysical modeling, enforcing physical and biological rules on machine learning approaches, uncertainty quantification, and time series data analysis. And I am looking for an opportunity to learn more about them.

% To sum up, given my analytical and physical insight and my programming experiences, I believe I am an ideal candidate for this Ph.D. program. I am eager and able to learn any kind of required prerequisites, languages, and courses to develop my own research project. I hope to do my Ph.D. in Computational Cancer Genomics/ Bioinformatics in your group and later become a professor engaged in research.

I am writing to express my strong interest in the PhD position in `Computational Cancer Genomics/Bioinformatics' in your group 
at the St. Anna Children’s Cancer Research Institute. 
With a master’s degree in Numerical Analysis and a solid background in applied mathematics, data analysis, and machine learning, 
I am eager to contribute to your research on the evolution and mutational processes of pediatric cancers.

I completed my master’s degree in Numerical Analysis at the Institute for Advanced Studies in Basic Sciences, 
where I developed a strong foundation in biophysical and mathematical modeling. 
Through the joint winter school with Graz University (as mentioned in my CV) and my own research on integrating imaging techniques to construct 3D breast cancer models, 
I became deeply interested in applying theoretical models to clinical problems. 
My master’s thesis focused on employing neural networks to solve partial differential equations, 
highlighting my interest in bridging applied mathematics with practical medical applications. 
Later, I directed my attention toward cancer modeling, 
studying works such as “Spatio-genetic and phenotypic modelling elucidates resistance and re-sensitisation to treatment in heterogeneous melanoma”, 
which presents a cutting-edge mathematical framework for investigating drug resistance. 
Its focus on mutational processes and tumor evolution strongly resonated with the themes emphasized in your group’s research, 
particularly in understanding which children are most at risk of poor treatment response.

As I explored your research, I was particularly drawn to your papers “Model-based tumor subclonal reconstruction” and 
“Exploiting evolutionary steering to induce collateral drug sensitivity in cancer”. 
I find your application of Physics-Informed Neural Networks (PINNs) especially compelling, 
as they achieve more reliable and interpretable results compared to purely data-driven approaches. 
I am also deeply impressed by how you integrate tumor evolution into your models and rigorously quantify data uncertainty. 
Moreover, I greatly admire the experimental design in your studies on drug resistance,
especially your consideration of both surviving and deceased cell populations and your hypothesis connecting these dynamics to the evolutionary history of the tumor. 
This comprehensive and mechanistic approach aligns closely with my own research philosophy.

I believe my background in applied mathematics, data analysis, and machine learning (particularly in modeling cancer biology) makes me a strong candidate for this position. 
My professional experience as a software developer has strengthened my programming, teamwork, and problem-solving skills, and 
taught me to perform effectively under pressure and tight deadlines. 
Beyond the technical aspects, my motivation to study cancer biology is deeply personal, having lost my aunt to breast cancer. 
I am eager to translate my skills and passion into developing computational methods that predict treatment response and improve patient outcomes. 
My research interests include oncology, biophysical modeling, physics-informed machine learning, uncertainty quantification, and time-series data analysis. 
Ultimately, I aspire to pursue an academic career where I can combine computational modeling with biological insight to advance cancer research.

Thank you very much for considering my application. I would be honored to contribute to your group and further discuss how my background and interests align with your research goals.

Sincerely,

\vspace{1em}
Sajed Zarinpour Nashroudkoli

\clearpage
